\hypertarget{appendix-7f1e17dcafb1}{}
\section*{Robust norm and intermediate range}
\section{Robust norm and intermediate range}\label{robust-norm-and-intermediate-range}

\textbf{Curriculum depth:} Introduced · \textbf{First seen:} C004 · \textbf{Latest development:} C004

\subsection{Definition and mechanism}\label{definition-and-mechanism}

A final representable answer can be lost to overflowing intermediates. Squaring 1e20f overflows before a finite norm could be formed. Scale by maximum component or use hypot to avoid unnecessary squared overflow.

\subsection{Prerequisites}\label{prerequisites}

\begin{itemize}
\tightlist
\item
  \hyperref[appendix-343615f49ddf]{IEEE-754 floating point} --- The format/arithmetic contract gives these bits their numerical domain.
\item
  \hyperref[appendix-eb336149c00a]{Floating-point infinities} --- Non-finite operands and overflowing intermediates need branches before ordinary finite formulas.
\end{itemize}

\subsection{Used by}\label{used-by}

\begin{itemize}
\tightlist
\item
  \href{}{Unit directions and safe normalization (02-Foundations-II, appendix-f744e41b95b4)} --- Phase 6 uses this established foundation in its geometry narrative.
\end{itemize}

Connect later dependent concepts here as their lessons arrive.

\subsection{Recall and next depth}\label{recall-and-next-depth}

Reconstruct the concrete example in the definition and identify its validity assumptions.

\textbf{Pending:} Non-finite policy, final range, accuracy and performance remain distinct; no repository fix is authorized.

\subsection{Source and navigation}\label{source-and-navigation}

\href{Continuous\%20Notes/Computing\%20Foundations\%20-\%20From\%20Bits\%20to\%20C++\%20and\%20Systems\%20APIs\#A\%20stretching\%20ruler\%20for\%20real-valued\%20quantities}{M004 teaching chapter} · \hyperref[appendix-fd090d982a15]{Per-part definitions} · \href{}{Dependency Map (02-Foundations-II, appendix-4ad47bdd269a)}
